\documentclass[11pt]{article}
\usepackage[utf8]{inputenc}
\usepackage[T1]{fontenc}
\usepackage{lmodern}
\usepackage[a4paper,margin=2.3cm]{geometry}
\usepackage{booktabs}
\usepackage{array}
\usepackage{xcolor}
\usepackage{enumitem}
\usepackage{parskip}
\usepackage{fancyhdr}
\usepackage[hidelinks]{hyperref}

\definecolor{accent}{HTML}{E2562B}
\newcommand{\code}[1]{\texttt{#1}}
\setlist[itemize]{leftmargin=1.4em}
\setlength{\emergencystretch}{2em}

\pagestyle{fancy}
\fancyhf{}
\lhead{\small\textcolor{accent}{AirStems}}
\rhead{\small Report de funcionalidad}
\rfoot{\small\thepage}
\renewcommand{\headrulewidth}{0.4pt}

\newcommand{\statusbox}[1]{%
  \noindent\fcolorbox{black!25}{black!4}{%
  \parbox{\dimexpr\linewidth-2\fboxsep-2\fboxrule}{#1}}\par\vspace{1em}}

\begin{document}

\begin{center}
  {\LARGE\bfseries\textcolor{accent}{AirStems}}\\[3pt]
  {\large Report de funcionalidad}\\[6pt]
  {\small Lluís Estapé \quad$\cdot$\quad Musicathon by Musixmatch Pro 2026 \quad$\cdot$\quad 2026-06-15}
\end{center}
\vspace{0.4em}
\hrule
\vspace{1.2em}

\noindent\textbf{Qué es.} Un instrumento que convierte la \emph{separación de stems} y la
\emph{letra sincronizada} en una \textbf{performance en directo}. Una webcam te sigue las manos
y, con gestos, remezclas en tiempo real las pistas de una canción (voz / batería / bajo / otros),
con los cambios \textbf{cuadrados al beat} y la \textbf{letra en karaoke}.

\vspace{1em}
\statusbox{\textbf{Estado (2026-06-15).} Núcleo \textbf{verificado offline}: separación local
(Demucs), motor de audio \textbf{estéreo}, detección de beats sobre el stem de batería
(129.2 BPM / 363 beats en la canción de prueba) y la app con tracking de manos + HUD + letra.
\textbf{APIs:} Musixmatch (letra sincronizada \emph{y} rich sync palabra-por-palabra)
\textbf{verificada}; LALAL.AI \emph{upload} verificado (la separación requiere la licencia premium
del hackathon); Cyanite (BPM/tono/mood) \textbf{verificada} de extremo a extremo (con MP3; la API falla en WAV).}

\tableofcontents
\vspace{1em}

\section{Cómo funciona: dos flujos}

\textbf{Flujo de AUDIO} (qué suena):
{\footnotesize\begin{verbatim}
  Cancion
    |
    v
  Separador  (Demucs  o  LALAL.AI)
    |
    v
  stems/<cancion>/{vocals,drums,bass,other}.wav
    |
    v
  StemEngine  ->  carga stems + detecta beat grid (librosa)
    |   (callback sounddevice, 44.1 kHz, bloque 512)
    v
  mezcla (gain por stem) -> filtro paso-bajo -> reverb -> tremolo -> SALIDA
\end{verbatim}}

\textbf{Flujo de CONTROL} (qué haces tú):
{\footnotesize\begin{verbatim}
  Webcam -> MediaPipe (2 manos) -> landmarks
    |
    |-- mano derecha  : dedos arriba/abajo    -> gain ON/OFF por stem (histeresis)
    |-- mano izquierda: altura de la muneca   -> filtro paso-bajo
    |                   abrir / cerrar la mano  -> reverb
    |-- beat-sync (tecla b)                   -> los cambios de gain saltan al siguiente beat
    |
  Letra (.lrc por linea  o  rich sync palabra-por-palabra) -> karaoke sincronizado
\end{verbatim}}

\section{Controles}

Primero \textbf{haz clic en la ventana ``AirStems''} para que tenga el foco del teclado.

\begin{center}
\begin{tabular}{@{}l p{0.60\linewidth}@{}}
\toprule
\textbf{Acción} & \textbf{Gesto / tecla} \\
\midrule
Play / pausa & \code{espacio} \\
Beat-sync ON/OFF & \code{b} \\
Cambiar de canción (siguiente) & \code{n} \\
Info / panel de controles & \code{i} \\
Salir & \code{q} \\
Stems 1--4 ON/OFF & mano derecha: subir/bajar cada dedo (índice$\to$meñique) \\
Drop total / mezcla completa & puño / mano abierta \\
Filtro paso-bajo & mano izquierda: altura de la muñeca (abajo oscuro, arriba abierto) \\
Reverb & mano izquierda: abrir / cerrar la mano (abierta full, puño 0) \\
\bottomrule
\end{tabular}
\end{center}

Los dedos usan \textbf{histéresis} (\code{UP\_MARGIN}/\code{DOWN\_MARGIN}) para no parpadear en el
umbral; el filtro y la reverb usan \textbf{suavizado EMA} (\code{SMOOTH}).

\section{El diferenciador: toggles cuadrados al beat}

Al cargar la canción, \textbf{librosa} detecta el BPM y el tiempo de cada beat. Con beat-sync
\textbf{ON}, cuando muteas o activas un stem el cambio \textbf{no es instantáneo}: se guarda como
\emph{pendiente} y se aplica en el \textbf{siguiente beat} (el callback de audio comprueba si cae un
beat dentro del bloque actual; el \emph{wrap} del bucle cuenta como \emph{downbeat}). Así cada
\emph{drop} y cada entrada \textbf{caen a tiempo} y suenan intencionados. Con beat-sync
\textbf{OFF}, los cambios son inmediatos. Las transiciones de gain usan una \textbf{rampa por
bloque}, así que no hay clics.

\section{Motor de audio (\texttt{StemEngine})}

\begin{itemize}
  \item Reproduce en bucle N stems (\textbf{estéreo}, 44.1 kHz) en un callback de \textbf{sounddevice} (bloque 512).
  \item \textbf{Mezcla} $=$ suma de (stem $\times$ gain) con rampa de gain por bloque, $\times$ master; efectos \textbf{por canal}.
  \item \textbf{Cadena de efectos} (reutilizada de \code{synth.py} de \emph{Aetheric Geometry}):
  \begin{itemize}
    \item Filtro \textbf{IIR paso-bajo} de 1.er orden (corte 200--8000 Hz según la altura de la mano).
    \item \textbf{Reverb Schroeder} (4 comb filters, feedback 0.82), \emph{wet} suavizado.
    \item \textbf{Trémolo} (LFO 0.5--10 Hz).
  \end{itemize}
  \item Expone: \code{position\_seconds} (sincroniza la letra), \code{bpm}, \code{beat\_pulse} (HUD) y \code{quantize}.
\end{itemize}

\section{Letra / karaoke (\texttt{lyrics.py})}

Parsea un \code{.lrc} (timestamps \code{[mm:ss.xx]}) y devuelve la \textbf{línea activa} según la
posición de reproducción del motor. Además soporta el \textbf{rich sync} de Musixmatch
(\code{.richsync.json}): timing \textbf{palabra por palabra}, así que el HUD resalta la parte ya
cantada (\code{current\_karaoke}). Si hay rich sync se prefiere; si no, cae al \code{.lrc}.

\section{Origen de los stems (desacoplado) y APIs de partners}

El motor \textbf{solo carga WAVs} de \code{stems/<cancion>/}: le da igual quién los generó. Por eso
las fuentes son \textbf{intercambiables}. Además descubre \textbf{todas} las canciones bajo
\code{stems/} y se alternan \textbf{en directo} con \code{n}, emparejando \code{lyrics/<nombre>}
y \code{analysis/<nombre>} por nombre.

\begin{center}
\begin{tabular}{@{}l l p{0.48\linewidth}@{}}
\toprule
\textbf{Módulo} & \textbf{Partner} & \textbf{Rol} \\
\midrule
\code{separate\_local.py} & --- (Demucs) & separación local/offline, sin key (desarrollo y \emph{fallback}) \\
\code{lalalai.py} & LALAL.AI & separación en la nube (v1 API), integración oficial \\
\code{musixmatch.py} & Musixmatch & descarga la letra sincronizada (\code{.lrc}) \\
\code{cyanite.py} & Cyanite & BPM / tono / mood (librosa hace la rejilla de beats local) \\
\bottomrule
\end{tabular}
\end{center}

\section{HUD (lo que se ve en pantalla)}
\begin{itemize}
  \item Barra por stem (verde $=$ ON, gris $=$ off).
  \item Barras de \textbf{filtro} y \textbf{reverb}.
  \item \textbf{Pulso} circular en cada beat $+$ \textbf{BPM} $+$ estado de beat-sync.
  \item \textbf{Línea de letra} sincronizada.
  \item Contador de \textbf{FPS}.
\end{itemize}

\section{Componentes (ficheros)}

\begin{center}
\begin{tabular}{@{}l p{0.52\linewidth}@{}}
\toprule
\textbf{Fichero} & \textbf{Rol} \\
\midrule
\code{airstems.py} & bucle principal: cámara, tracking, mapeo gestos$\to$params, HUD, teclas \\
\code{stem\_engine.py} & \code{StemEngine}: audio en tiempo real, mezcla, efectos, beat grid + beat-sync \\
\code{lyrics.py} & parser LRC + línea actual \\
\code{separate\_local.py} & separación local con Demucs (carga audio vía soundfile) \\
\code{lalalai.py} & cliente LALAL.AI v1 (upload $\to$ split $\to$ check $\to$ download) \\
\code{musixmatch.py} & cliente de letra sincronizada \\
\code{cyanite.py} & cliente de análisis (BPM/tono/mood) \\
\code{smoke\_test.py} & verificación headless (sin cámara ni altavoz) \\
\code{gestures.py}, \code{renderer.py}, \code{config.py} & reutilizados de \emph{Aetheric Geometry} \\
\bottomrule
\end{tabular}
\end{center}

\section{Stack técnico}

Python $\cdot$ OpenCV $\cdot$ MediaPipe $\cdot$ sounddevice $\cdot$ soundfile $\cdot$ NumPy $\cdot$
librosa $\cdot$ (Demucs / torch para la separación local) $\cdot$ APIs: \textbf{LALAL.AI},
\textbf{Musixmatch}, \textbf{Cyanite}.

\section{Cómo ejecutar}

Con el venv de \emph{Aetheric} (mediapipe 0.10.9, que sí tiene \code{mp.solutions}):
{\footnotesize\begin{verbatim}
& "...\Gestural-Harmonic-Mapping\venv\Scripts\python.exe" "...\AirStems\airstems.py"
\end{verbatim}}
Separar otra canción:
{\footnotesize\begin{verbatim}
python separate_local.py "C:\ruta\cancion.wav"   # local (Demucs)
python lalalai.py        "C:\ruta\cancion.wav"   # LALAL.AI (con key)
\end{verbatim}}

\section{Pendiente / próximos pasos}
\begin{itemize}
  \item \textbf{Cyanite:} \textbf{verificado} de extremo a extremo --- análisis V7 (BPM/tono/mood)
        funcionando con MP3 (la API falla en WAV PCM-16). Falta surtir BPM/mood al HUD.
  \item \textbf{LALAL.AI:} correr una separación real al activarse la licencia premium del hackathon
        (hoy: \emph{upload} verificado; Demucs cubre desarrollo y demo).
  \item \textbf{Mío (post-examen):} afinar umbrales de gestos en vivo, elegir tema y grabar el vídeo.
  \item \textbf{Stretch:} export de la remezcla a WAV $\cdot$ \emph{pinch} $=$ solo de un stem $\cdot$
        volumen master con la distancia entre manos.
\end{itemize}

\end{document}
